\section{The TIGER Dimensions}
\label{sec:method}

The five TIGER dimensions are derived from the top PCs of the within-game normalized
TIGRIS corpus (see TG-A2 for the full derivation). Each dimension is formalized here
with a name, a one-sentence definition, and anchor descriptors at 0, 5, and 10 on
a normalized scale.

\subsection{Tension (T)}

\textbf{Definition:} The degree to which the game sustains meaningful decision pressure
through resource scarcity, timing constraints, or competitive threat.
\textit{Score 0:} Decisions are unconstrained; no meaningful pressure operates at any
point. \textit{Score 5:} Moderate pressure punctuates key decision nodes but relaxes
between them. \textit{Score 10:} Sustained high-pressure decision-making throughout;
the game is primarily an exercise in constraint navigation.

\subsection{Interaction (I)}

\textbf{Definition:} The degree to which player outcomes depend on direct engagement
with other players' choices, boards, or resources.

\subsection{Game-Architecture (G)}

\textbf{Definition:} The degree to which the game's structural design --- its sequence
of phases, its scoring timing, its feedback loops --- shapes strategic possibility
independently of moment-to-moment decisions.

\subsection{Elegance (E)}

\textbf{Definition:} The degree to which the game achieves deep strategic variety from
a small number of rules, with low surface complexity relative to strategic depth.

\subsection{Range (R)}

\textbf{Definition:} The degree to which the game supports a wide range of player types,
experience levels, and group compositions through its mechanical flexibility.
